\subsection{Invariant Factors}

\begin{definition} \label{def:operator-invariant-factors-elementary-divisors} \hlt{Invariant Factors and Elementary Divisors}\\
Let $T\in\LL(\mathcal{V})$.
\begin{enumerate}[label=(\roman*)]
\setlength{\itemsep}{0pt}
\item A \hlt{cyclic decomposition} is $\mathcal{V}=\bigoplus_{j=1}^r\mathcal{C}_j$, where each $\mathcal{C}_j$ is nonzero, $T$-invariant, and cyclic.
\item $\operatorname{InvFact}(T)\defeq(f_1,\ldots,f_r)$, where $f_j$ are monic nonconstant invariant factors of $\mathcal{V}_T$; $f_r\mid f_{r-1}\mid\cdots\mid f_1$.
\item $\operatorname{ElemDiv}(T)$ is the multiset of monic prime powers in the $f_j$, counted with multiplicity.
\item \hlt{Elementary divisor basis} is $\mathcal{B}_{\operatorname{ED}}\defeq\bigsqcup_{\alpha}(\mathbf{v}_{\alpha},T\mathbf{v}_{\alpha},\ldots,T^{d_{\alpha}-1}\mathbf{v}_{\alpha})$, where $d_{\alpha}=\deg m_{\mathbf{v}_{\alpha}}$; $\{m_{\mathbf{v}_{\alpha}}\}_{\alpha}=\operatorname{ElemDiv}(T)$.
\item For $A\in\operatorname{Mat}(n,n,\mathbb{F})$, set $\operatorname{InvFact}(A)\defeq\operatorname{InvFact}(T_A)$ and $\operatorname{ElemDiv}(A)\defeq\operatorname{ElemDiv}(T_A)$.
\end{enumerate}
\end{definition}

\begin{theorem} \label{thm:operator-primary-cyclic-decomposition} \hlt{Primary Cyclic Decomposition for an Operator}\\
Let $\mu_T=\prod_{i=1}^sp_i^{e_i}$ with $p_i\in\mathbb{F}[z]$ distinct monic irreducible, and set $\mathcal{V}_{p_i}\defeq\ker p_i(T)^{e_i}$. Then
\begin{enumerate}[label=(\roman*)]
\setlength{\itemsep}{0pt}
\item $\mathcal{V}_T=\bigoplus_{i=1}^s\mathcal{V}_{p_i}$;
\item $\mathcal{V}_{p_i}=\bigoplus_{j=1}^{k_i}\mathcal{C}_T(\mathbf{v}_{i,j})$, with $m_{\mathbf{v}_{i,j}}=p_i^{e_{i,j}}$ and $e_i=e_{i,1}\geq e_{i,2}\geq\cdots\geq e_{i,k_i}$;
\item $\operatorname{ElemDiv}(T)=\{p_i^{e_{i,j}}:1\leq i\leq s,\ 1\leq j\leq k_i\}$, corresponding cyclic bases form elementary divisor basis.
\end{enumerate}
\begin{proof}
\begin{enumerate}[label=(\roman*)]
\setlength{\itemsep}{0pt}
\item By Theorem \ref{thm:primary-cyclic-decomposition-pid}, $\mathcal{V}_T$, $\mathcal{V}_T=\bigoplus_i\mathcal{V}_T[p_i^\infty]$; $\mathcal{V}_T[p_i^\infty]\defeq\{\mathbf{x}\in\mathcal{V}:\exists N\geq1,\ p_i(T)^N\mathbf{x}=0\}=\ker p_i(T)^{e_i}=\mathcal{V}_{p_i}$.
\item By Theorem \ref{thm:primary-cyclic-decomposition-pid}, $\mathcal{V}_{p_i}=\bigoplus_{j=1}^{k_i}\mathcal{C}_T(\mathbf{v}_{i,j})$, $m_{\mathbf{v}_{i,j}}=p_i^{e_{i,j}}$ and $e_i=e_{i,1}\geq\cdots\geq e_{i,k_i}$.
\item Definition \ref{def:operator-invariant-factors-elementary-divisors} gives $\operatorname{ElemDiv}(T)=\{p_i^{e_{i,j}}:1\leq i\leq s,\ 1\leq j\leq k_i\}$, and concatenating the bases of $\mathcal{C}_T(\mathbf{v}_{i,j})$ gives an elementary divisor basis.
\end{enumerate}
\hlt{Structural Commentary}: Primary module pieces are polynomial-kernel pieces for the operator.\\
\hlt{Key Insight}: Elementary divisors are exactly the orders of the cyclic primary summands.
\end{proof}
\end{theorem}

\begin{theorem} \label{thm:operator-invariant-factor-skeleton} \hlt{Invariant-Factor Skeleton for Operators}\\
Let $0<\dim\mathcal{V}<\infty$, $T\in\LL(\mathcal{V})$, and $\operatorname{InvFact}(T)=(f_1,\ldots,f_r)$.\\
Then $f_r\mid\cdots\mid f_1$ and there are $T$-cyclic subspaces $\mathcal{C}_j$ such that $\mathcal{V}=\bigoplus_{j=1}^r\mathcal{C}_j$, $\mathcal{C}_j\cong\mathbb{F}[z]/\langle f_j\rangle$, $[T]=\bigoplus_{j=1}^r C_{f_j}$ in the concatenated cyclic bases, $\mu_T=f_1$, and $p_T=\prod_{j=1}^r f_j$.
\begin{proof}
By Proposition \ref{prop:operator-modules-are-fg-torsion}, $\mathcal{V}_T$ is finitely generated torsion over $\mathbb{F}[z]$.\\
Theorem \ref{thm:pid-invariant-factor-decomposition} gives $\mathcal{V}_T\cong\bigoplus_{j=1}^r\mathbb{F}[z]/\langle f_j\rangle$ with $f_r\mid\cdots\mid f_1$, uniquely.\\
Transport cyclic generators to $T$-cyclic subspaces $\mathcal{C}_j$; Theorem \ref{thm:cyclic-subspace-quotient-and-companion} gives $\mathcal{C}_j\cong\mathbb{F}[z]/\langle f_j\rangle$ and block $C_{f_j}$.\\
Since $\operatorname{Ann}(\mathcal{V}_T)=\langle\operatorname{lcm}(f_1,\ldots,f_r)\rangle=\langle f_1\rangle=\langle\mu_T\rangle$, $f_1=\mu_T$.\\
Finally, Theorem \ref{thm:companion-matrix-has-prescribed-minimal-polynomial} gives $p_{C_{f_j}}=f_j$, so $p_T=\prod_{j=1}^r f_j$.\\
\hlt{Structural Commentary}: Torsion module decomposition produces companion blocks.\\
\hlt{Key Insight}: Invariant factors are the polynomial-module version of choosing canonical blocks for similarity.
\end{proof}
\end{theorem}
